%قبل از شروع مسئله ترجیحاً چیزی ننویسید. اما اگر خواستید بنویسید، در قالب لیستی (فقرات و شمارش) بنویسید.
%برای اضافه کردن نمره، بعدی از شروع‌مسئله، [x] را اضافه کنید و به جای x نمره مورد نظر را قرار دهید. برای زیرمسئله‌ها هم به همین صورت است.
%\شروع‌مسئله[10]
\شروع‌مسئله 
در این چالش شما مرحله به مرحله یک الگوریتم تشخیص برخورد دوبعدی را پیاده‌سازی می‌کنید. این مسأله در ردۀ هندسۀ محاسباتی قرار می‌گیرد. در مسألۀ تشخیص برخورد دوبعدی تعدادی چندضلعی (که لزوماً محدب نیستند) به شما داده می‌شود و تمام زوج‌هایی که برخورد دارند از شما خواسته می‌شود. دقت کنید که مماس بودن دو چندضلعی به معنای برخورد آن‌ها نیست و زمانی دو چندضلعی برخورد دارند که مساحت ناحیۀ مشترک آن‌ها مثبت باشد. این مسأله در بازی‌های کامپیوتری، شبیه‌سازی‌های فیزیکی و روباتیک کاربرد دارد. بسیاری از کسانی که به ساخت بازی‌های کامپیوتری می‌پردازند به طور مستقیم و یا غیرمستقیم از الگوریتم‌های تشخیص برخورد استفاده می‌کنند. هدف از این چالش، آشنایی شما با الگوریتم‌های تشخیص برخورد و استفاده از داده‌ساختارهای مناسب برای بهینه کردن عملیات تشخیص برخورد است. شما در سه مرحله یک الگوریتم تشخیص برخورد پیاده‌سازی می‌کنید که به شرح زیر است:
\begin{itemize}
	\item \textbf{مرحلۀ اول:}
	در این مرحله شما الگوریتمی را پیاده‌سازی می‌کنید که برخورد و یا عدم برخورد بین دو چندضلعی را تشخیص می‌دهد.
	\item \textbf{مرحلۀ دوم:}
	در این مرحله شما یک الگوریتم \lr{space partitioning} را پیاده‌سازی خواهید کرد که در آینده با این نوع الگوریتم‌ها آشنا خواهید شد.
	\item \textbf{مرحلۀ سوم:}
	در مرحلۀ آخر با استفاده از کدهایی که در مراحل قبل زدید یک الگوریتم تشخیص برخورد را به صورت کامل پیاده‌سازی خواهید کرد. در این مرحله می‌توانید کدتان را تا جای ممکن بهینه کنید.
\end{itemize}
برای آشنایی بیشتر با این مسأله می‌توانید 
\href{https://en.wikipedia.org/wiki/Collision_detection}{این لینک}
را مطالعه کنید. برای آشنایی بیشتر با مراحل بهینه‌سازی نیز می‌توانید بخش \lr{optimization} لینک داده شده را مطالعه کنید. جزئیات هر مرحله متعاقباً اعلام خواهد شد.

\متن‌سیاه{نحوۀ نمره‌دهی:}

مرحلۀ اول و دوم هر کدام حداکثر نیم نمره و مرحلۀ سوم حداکثر یک نمره دارد. نمره‌دهی در مرحلۀ اول و دوم توسط کوئرا انجام می‌شود. به این صورت که فرد بر اساس تعداد تست کیس‌هایی که درست پاسخ می‌دهد نمره می‌گیرد. نمره‌دهی در مرحلۀ آخر به صورت رقابتی است. به این معنا که مجموع زمان هر نفر بر روی تست کیس‌ها محاسبه می‌شود و افراد بر اساس این مجموع زمان مرتب می‌شوند. سپس به نفر اول یک نمره تعلق می‌گیرد و نمره بقیۀ افراد به صورت خطی کاهش می‌یابد.

%\متن‌سیاه{سوال عملی اول:}\\
%\متن‌سیاه{دراپ‌تست}\\

%\متن‌سیاه{ورودی:}\\
%در خط اول ورودی، ابتدا k و سپس m داده می‌شود.\\
%\متن‌سیاه{خروجی:}\\
%در تنها سطر خروجی، تعداد آیفون‌های مورد نیاز می‌آید. اگر در هیچ صورتی چنین کاری ممکن نبود ۱- چاپ می‌شود.\\
%\متن‌سیاه{محدودیت‌ها:}\\
%$1\leq k\leq 10^4$

%$m\leq10^9$
%\\
%\textbf{\متن‌سیاه{ورودی و خروجی نمونه}\\
%% تست اول
%\begin{latin}
% \begin{tabularx}{\textwidth}{|X | X|}
% \hline
% \textbf{stdin} &
% \textbf{stdout}\\
% \hline
% 100 100&
% 1\\
% \hline
% 
% \end{tabularx}
%\end{latin}
%
%% تست دوم
%
%\begin{latin}
% \begin{tabularx}{\textwidth}{|X | X|}
% \hline
% \textbf{stdin} &
% \textbf{stdout}\\
% \hline
% 340 20&
% 3\\
% \hline
% 
% \end{tabularx}
%\end{latin}}
% صورت اصلی مسئله را در این‌جا بنویسید.
%اگر مسئله چند قسمت داشت به این صورت عمل کنید:
%\شروع{زیرمسائل}
%\زیرمسئله[0] % مسئلۀ ۱آ
%\زیرمسئله % مسئلۀ ۱ب
%\زیرمسئله % مسئلۀ ۱پ
%\زیرمسئله % مسئلۀ ۱ت
%\پایان{زیرمسائل}

\solution{
	 1)
	$$O({n \choose 4}) < O(n^5) < O(2^n) < O(n!) $$
	2)
	$$O(n\log n ) <O(n\sqrt{n}) <O((\log n) ^ {100}) <O(2 ^{ \sqrt{n}}) $$
	3)
	$$ O(\sum_{i=1}^{n} \frac{1}{i}) = O(\log n) \Rightarrow$$
	$$O(1000! ) <O(n \sum_{i=1}^{n} \frac{1}{i}) <O(n^2 \log(\log n)) <O(2 ^ {\log n}) $$
پاسخ نمونه.
}




%%%%%%%%%%%%%%%%%%%%%%%%%%%%%%%%%%%%%%%%%%%%%%%%%%%%%
% برخی دستورات مفید برای تمرین‌ها. %
%%%%%%%%%%%%%%%%%%%%%%%%%%%%%%%%%%%%%%%%%%%%%%%%%%%%%

% فرمول‌های نمونه:
%$f(x) = \Theta(x)$
%$T(x, y) \leq \log(x) + 2^y + \binom{2n}{n}$
% {\tt code\_function}


% لیست موارد بدون شماره‌گذاری:
%\شروع{فقرات}
% \فقره مورد اول 
% \شروع{فقرات}
% \فقره مورد اول داخلی 
% \شروع{فقرات}
% \فقره اولین مورد خیلی داخلی 
% \فقره دومین مورد خیلی داخلی 
% \پایان{فقرات}
% \فقره مورد دوم داخلی 
% \پایان{فقرات}
% \فقره مورد دوم 
%\پایان{فقرات}

% لیست موارد شماره‌گذاری شده:
%\شروع{شمارش}
% \فقره یک مورد 
% \فقره یک مورد دیگر 
% \فقره باز هم یک مورد دیگر
%\پایان{شمارش}

% برای نوشتن یک فرمول چند خطی که یک شماره بخوره. برای حالت‌های دیگر به سایت زیر مراجعه کنید:
% https://www.overleaf.com/learn/latex/Aligning_equations_with_amsmath
%\begin{align} 
% \begin{split}
% (x+y)^3 &= (x+y)^2(x+y) \\
% &= (x^2+2xy+y^2)(x+y) \\
% &= (x^3+2x^2y+xy^2) + (x^2y+2xy^2+y^3) \\
% &= x^3+3x^2y+3xy^2+y^3
% \end{split} 
%\end{align}

% نوشتن ماتریس و آرایه دوبعدی
%\begin{align}
% A = 
% \begin{bmatrix}
% A_{11} & A_{21} \\
% A_{21} & A_{22}
% \end{bmatrix}
%\end{align}

% اضافه کردن تصویر با شرح. برای اضافه کردن تصویر بدون شرح، فقط دستور درج‌تصویر را استفاده کنید.
%\شروع{شکل}[h!]
%\تنظیم‌ازوسط
% \درج‌تصویر[پهنا=0.5\پهنای‌متن]{img.jpg}
%\شرح{شرح تصویر}
%\برچسب{تصویر:۱}
%\پایان{شکل}
%\پایان{زیرمسائل}

%کد زیر را ببینید.
%%\lstinputlisting{source_filename.py}
%\begin{latin}
%\begin{lstlisting}
%# Import the modules
%import sys
%import random
%
%ans = True
%
%while ans:
%question = raw_input("Ask the magic 8 ball a question: (press enter to quit) ")
%
%answers = random.randint(1,8)
%
%if question == "":
%sys.exit()
%
%elif answers == 1:
%print "It is certain"
%
%elif answers == 2:
%print "Outlook good"
%
%elif answers == 3:
%print "You may rely on it"
%
%elif answers == 4:
%print "Ask again later"
%
%elif answers == 5:
%print "Concentrate and ask again"
%
%elif answers == 6:
%print "Reply hazy, try again"
%
%elif answers == 7:
%print "My reply is no"
%
%elif answers == 8:
%print "My sources say no"
%\end{lstlisting}
%\end{latin}